% Section 5.9, from the summary table of lectures/L05/appendix/a_theory.md and from
% lectures/L05/README.md: the objectives, the evaluation questions and the next lesson.

\section{Sammanfattning}
\label{c5:sec:sammanfattning}

\begin{booktable}{@{}p{12.5em}L@{}}
  \thead{Begrepp} & \thead{Nyckelformel} \\
  \midrule
  Produktregeln & $a^m \cdot a^n = a^{m+n}$ \\
  Kvotregeln & $a^m / a^n = a^{m-n}$ \\
  Potens av potens & $(a^m)^n = a^{mn}$ \\
  Negativ exponent & $a^{-n} = 1/a^n$ \\
  Kvadratrot & $\sqrt{a} = a^{1/2}$ \\
  Bråkexponent & $a^{m/n} = \sqrt[n]{a^m}$ \\
  Standardform & $a \times 10^n$, $1 \leq a < 10$ \\
  Effekt & $P = U^2/R = RI^2$ \\
  Positionssystem & $1011_2 = 1 \cdot 2^3 + 0 \cdot 2^2 + 1 \cdot 2^1 + 1 \cdot 2^0 = 11$ \\
  Binärt och hexadecimalt & En hexadecimal siffra motsvarar fyra bitar \\
  Antal värden med $n$ bitar & $2^n$, från $0$ till $2^n - 1$ \\
\end{booktable}

\needspace{6\baselineskip}
\subsection*{Det här ska du kunna}
Efter det här kapitlet ska du:
\begin{itemize}
  \item Känna till och kunna tillämpa potensreglerna.
  \item Kunna beräkna och förenkla uttryck med rötter.
  \item Kunna skriva tal på standardform och hantera värdesiffror.
  \item Kunna omvandla tal mellan decimal, binär och hexadecimal form.
\end{itemize}

\testadigsjalv
\begin{enumerate}
  \item Förenkla $\dfrac{(3x^2)^3}{9x^4}$ med potensreglerna.
  \item En kapacitans är $C = 220\,\text{nF}$. Ange värdet i farad på standardform.
  \item Omvandla $1101\,0010_2$ till hexadecimal och decimal form.
\end{enumerate}

\needspace{6\baselineskip}
\subsection*{Nästa kapitel}
\Kapref{6} handlar om andragradsekvationer: pq-formeln och kvadratkomplettering.
